\abstract{%
  My abstract ...
}

% The abstract is a short, self-contained statement describing the whole of your work.  It should be less than a page (typically half a page) and should summarise the scope, purpose, results and content of the work.  The abstract might be thought of as a summary which you would read quickly to decide if the rest of the document is worth taking the time to read in detail.  In scientific publishing, abstracts are often used as sources of keywords and concepts for searching, so it’s important to ensure that the main ideas and conclusions of your work are present.

% When someone has read your abstract, they should know what your project was about, how you did it and what the end result was.  It doesn’t need to contain references or literature reviews.  For example, here’s the abstract from Bell and Brooks (2019)’ s paper “What makes students satisfied? A discussion and analysis of the UK's national student survey” (which, incidentally, is very interesting):

% “This paper analyses data from the National Students Survey, determining which groups of students expressed the greatest levels of satisfaction. We find students registered on clinical degrees and those studying humanities to be the most satisfied, with those in general engineering and media studies the least. We also find contentment to be higher among part-time students, and significantly higher among Russell group and post-1992 universities. We further investigate the sub-areas that drive overall student satisfaction, finding teaching and course organisation to be the most important aspects, with resources and assessment and feedback far less relevant. We then develop a multi-attribute measure of satisfaction which we argue produces a more accurate and more stable reflection of overall student satisfaction than that based on a single question.”